\chapter{Investigation/Experiment, Result, Analysis and Discussion}
\label{ch:results}

This chapter reports the evaluation of LectureAssist's question generator. It
describes the environment the experiment was run in, defines the evaluation
protocol precisely, and reports the results. The evaluation adopts the protocol
of \textbf{EQGBench}~\cite{eqgbench}, a published benchmark for evaluating an
LLM's ability to \emph{generate} educational questions, so that the four
generation pipelines are scored on the same criteria a purpose-built
question-generation benchmark uses, rather than on ad-hoc metrics.

The experiment is designed around three questions.
\emph{RQ1:} judged on the EQGBench dimensions, does the agentic pipeline
produce better multiple-choice questions than single-shot prompting, including
the stronger chain-of-thought and program-of-thought prompting strategies?
\emph{RQ2:} how does each pipeline behave as the target difficulty rises from
easy to hard?
\emph{RQ3:} which dimensions actually separate the pipelines, and which
saturate?

\section{Environment Setup}
\label{sec:res-env}

Question \emph{generation} was run on a single Google Colab instance with an
NVIDIA T4 GPU ($16$\,GB), with models served locally by Ollama; no paid or
proprietary API is used anywhere in the generation path. Bulk-generation
artifacts are written to Google Drive so a run survives a Colab disconnection
and can be resumed rather than restarted.

The generator is \texttt{gemma3:4b}, pinned: it is never allowed to escalate to
a larger model, even when it returns malformed output, so that ``generated by
the $4$B'' is true of every question in the study.

\paragraph{Pipelines.} Four generation pipelines are compared, all using the
same pinned generator and the same source material:
\begin{itemize}
  \item \textbf{Agentic} --- the production per-question loop of
        \emph{Plan $\rightarrow$ Generate $\rightarrow$ Validate $\rightarrow$
        Retry}: a planner selects topics, each question is checked by a
        self-validator on grounding, difficulty, instruction compliance and
        uniqueness, and regenerated (with the failure reason fed back) until it
        passes or a retry budget is exhausted; the generator may issue its own
        retrieval query when it judges the supplied context insufficient.
  \item \textbf{Non-agentic} --- a single LLM call that writes questions
        directly, with no planning, validation, retrieval or retry.
  \item \textbf{Chain-of-thought (CoT)} --- as non-agentic, but the prompt
        instructs the model to reason step-by-step about the content before
        writing each question.
  \item \textbf{Program-of-thought (PoT)} --- as non-agentic, but the model is
        instructed to derive each question through explicit worked computation
        before committing to an answer.
\end{itemize}
The three baseline pipelines read the document through an identical sliding
window that steps across the whole chapter with a uniform batch size, so no
baseline is advantaged by seeing more or different content, and any difference
between them reflects the prompting strategy alone.

\paragraph{Source data.} Questions are generated from a college-level STEM
document --- the \emph{Limits} chapter of OpenStax \emph{Calculus, Volume~1}
\cite{openstax_calculus} --- ingested through the standard document pipeline
(extraction $\rightarrow$ map-reduce summary $\rightarrow$ retrieval index). The
extracted chapter text (about $112$k characters) is also supplied to the judge as
the ground-truth reference against which each question is scored
(Section~\ref{sec:res-protocol}).

\paragraph{Question corpus.} Each pipeline generated $50$ MCQs at each of three
target difficulties (easy, medium, hard). After generation, questions that were
structurally malformed (an empty option, or no valid A--D answer key) were
removed, since a malformed item cannot be judged fairly on any dimension. The
evaluated corpus is summarised in Table~\ref{tab:corpus}; the per-cell counts
$N$ are carried into every results table because they are not uniform.

\begin{table}[H]
\centering
\caption{Evaluated question corpus after removing structurally malformed items.}
\label{tab:corpus}
\small
\begin{tabular}{|l|c|c|c|c|}
\hline
\textbf{Pipeline} & \textbf{Easy} & \textbf{Medium} & \textbf{Hard} & \textbf{Total} \\
\hline
Agentic     & 49 & 44 & 40 & 133 \\
Non-agentic & 50 & 50 & 50 & 150 \\
CoT         & 50 & 50 & 50 & 150 \\
PoT         & 50 & 50 & 50 & 150 \\
\hline
\textbf{Total} & 199 & 194 & 190 & 583 \\
\hline
\end{tabular}
\end{table}

\section{Evaluation Protocol}
\label{sec:res-protocol}

The evaluation follows EQGBench~\cite{eqgbench}. Generation and measurement are
kept separate: the Colab run only \emph{generates} and saves every question to a
single JSON artifact; a standalone harness then reads that artifact and scores
each question. Scoring an already-generated corpus with a separate program means
the measurement can be repeated or re-judged without touching the system under
test.

\paragraph{Dimensions and scale.} Following EQGBench, every question is scored on
five dimensions, each on a three-point scale $\{0,1,2\}$ (Poor / Good /
Excellent):
\begin{itemize}
  \item \textbf{KP --- Knowledge-Point Alignment:} does the question genuinely
        target the intended knowledge point, grounded in the source material?
  \item \textbf{QT --- Question-Type Alignment:} is it a well-formed item of the
        requested type (a four-option MCQ with exactly one defensible answer)?
  \item \textbf{QQ --- Question Item Quality:} is the stem clear and
        unambiguous, and are the options and distractors well constructed?
  \item \textbf{SQ --- Solution Explanation Quality:} is the provided
        explanation correct, complete and pedagogically useful?
  \item \textbf{CG --- Competence-Oriented Guidance:} does the item cultivate
        understanding, reasoning or application (higher-order competence) rather
        than rote recall?
\end{itemize}
The reported per-pipeline score on each dimension is the mean over that
pipeline's questions, and the \emph{composite} is the mean of the five
dimensions.

\paragraph{Judge and voting.} EQGBench scores each question with a
reasoning-model judge and reduces judge noise by voting over repeated rounds.
The same procedure is used here: every question is judged \textbf{three times
independently}, and the final score on each dimension is the \emph{mode} of the
three rounds (the median is used on the rare occasion all three disagree). Each
judge call \emph{fails closed} --- a call that errors or returns unparseable
output contributes a zero, never an inflated score --- and a question whose three
rounds all fail is scored zero on every dimension rather than dropped.

\paragraph{Judge models.} EQGBench uses DeepSeek-R1 as its judge; that model was
not reachable through an accessible endpoint when this evaluation was run, so it
is not used here. Instead every question is scored by \emph{two} judges and both
are reported. The first is \texttt{gemma3:12b}, the larger local model, run
through the same offline harness; note that although it is a bigger model than
the \texttt{gemma3:4b} generator, it belongs to the same \texttt{gemma} family,
so it may share some of the generator's blind spots. The second is
\texttt{gpt-oss-120b}, a $120$B open-weight reasoning model from a different
family entirely, which serves as the cross-family check: where a same-family and
a cross-family judge agree, a pipeline's ranking cannot easily be dismissed as an
artifact of one model's quirks. Using two judges of different lineage --- and
reporting how far they agree --- is how this evaluation answers the objection
that an LLM judging an LLM's output is merely grading itself. What is adopted
from EQGBench is the \emph{protocol} --- the five dimensions, the $0$--$2$ scale,
and three-round voting --- not the specific judge model; results should be read
as EQGBench-protocol scores under these two independent judges.

\paragraph{Why not overlap metrics.} Following EQGBench, no $n$-gram overlap
metric (BLEU/ROUGE) is reported. Generation here is open-ended over a whole
chapter, so there is no reference question set to overlap with, and similarity
to any particular phrasing would not be evidence of quality.

\paragraph{Judge prompt.} Each question is sent to the judge with the source
material and the rubric below (reconstructed from the EQGBench dimension
definitions, as the benchmark has no public code release):

\begin{verbatim}
You are an expert evaluator of educational assessment items. Evaluate
this ONE multiple-choice question along the five EQGBench dimensions.
Score each dimension 0 (Poor), 1 (Good), or 2 (Excellent). Be strict;
2 is reserved for genuinely excellent performance on that dimension.

Reference material (ground truth): [chapter content]
Intended knowledge point (topic): [topic]
Requested question type: multiple-choice (4 options A-D, one correct)
Requested difficulty: [easy|medium|hard]

Question item:
Q: [question]      Options: [A-D]
Marked correct answer: [letter]
Provided solution explanation: [explanation]

Dimensions:
1. KP (Knowledge-Point Alignment)
2. QT (Question-Type Alignment)
3. QQ (Question Item Quality)
4. SQ (Solution Explanation Quality)
5. CG (Competence-Oriented Guidance)

Return STRICT JSON only:
{"KP":0,"QT":0,"QQ":0,"SQ":0,"CG":0}
\end{verbatim}

\section{Results and Discussion}
\label{sec:res-results}

\subsection{What EQGBench does}
\label{sec:res-what}

EQGBench is a published benchmark that asks a different question from most
evaluations. Instead of testing whether a model can \emph{answer} questions, it
tests whether a model can \emph{write} good educational questions. The authors
collected a large set of school-level science and mathematics material, asked
several large language models to generate questions from it, and then had a
strong ``judge'' model read each generated question and rate it. Every question
is rated on the five qualities listed earlier --- knowledge-point alignment,
question-type alignment, question quality, solution quality, and
competence-oriented guidance --- and each quality is given one of three grades:
poor, good, or excellent (written as $0$, $1$, and $2$). To keep a single
judge's opinion from dominating, each question is judged three times and the
grade that appears most often is kept.

The headline finding of EQGBench is the useful part for this project. On four of
the five qualities almost every model scored close to excellent, so those four
qualities did not tell the models apart. The fifth quality --- competence-oriented
guidance, that is, whether a question actually makes a student think rather than
just recall a fact --- was where every model scored poorly, and that is the
quality that separated the good models from the weak ones. Table~\ref{tab:eqg-their}
shows this pattern with figures reported in the EQGBench paper.

\begin{table}[H]
\centering
\caption{The pattern reported by EQGBench: the first four qualities are almost
maxed out for every model, while competence-oriented guidance (CG) collapses.
Figures are as reported in the EQGBench paper (mathematics subset) and are shown
to illustrate the pattern; exact values should be taken from the paper itself.}
\label{tab:eqg-their}
\small
\begin{tabular}{|l|c|c|}
\hline
\textbf{Quality (score out of 2)} & \textbf{Typical model} & \textbf{Pattern} \\
\hline
Knowledge-point alignment (KP) & $\approx 1.9$--$2.0$ & near-maximum for all \\
Question-type alignment (QT)   & $\approx 1.9$--$2.0$ & near-maximum for all \\
Question quality (QQ)          & $\approx 1.9$--$2.0$ & near-maximum for all \\
Solution quality (SQ)          & $\approx 1.9$--$2.0$ & near-maximum for all \\
Competence guidance (CG)       & $\approx 0.2$        & collapses for all \\
\hline
\end{tabular}
\end{table}

\subsection{How this fits our project}
\label{sec:res-fit}

Our system also \emph{writes} questions, so EQGBench is a natural fit: it was
built to measure exactly the task our generator performs. We therefore score our
four pipelines on the same five qualities, on the same three-point scale, using
the same three-times-and-take-the-most-common judging. This lets us report our
generator against an external, published yardstick rather than a scale we
invented ourselves.

Two honest differences should be stated plainly. First, EQGBench uses a specific
judge model (DeepSeek-R1); that model was not reachable when we ran our
evaluation, so we grade with two other strong models instead --- the local
\texttt{gemma3:12b} and a different-family reasoning model,
\texttt{gpt-oss-120b} --- and report both. Grading with two judges, one from the
same family as our generator and one from a different family, lets us show that a
result holds up no matter which judge is asked, rather than resting on a single
model's opinion. What we borrow from EQGBench is its method --- the five
qualities, the grading scale, and the repeated judging --- not its exact judge.
Second, EQGBench works on school-level science and mathematics, whereas our
questions come from a university calculus chapter; the method carries over, but
the subject matter is our own.

\subsection{How our evaluation works}
\label{sec:res-how}

The evaluation runs in two clearly separated stages. In the first stage the
system only \emph{generates}: all four pipelines produce their questions and
everything is saved to a single file, with no scoring involved. In the second
stage a separate program reads that file and does the scoring. Keeping the two
apart means the scoring can be repeated, or redone with a different judge, without
ever re-running or disturbing the generator.

For each saved question a judge is given the original chapter text as the source
of truth, together with the question, its options, its marked answer and its
explanation. The judge returns a grade of $0$, $1$ or $2$ on each of the five
qualities. This is repeated three times per question and the most common grade is
kept. If a judging attempt fails or returns something unreadable it is treated as
a zero rather than quietly ignored, so a broken question can never accidentally
look good. Each pipeline's score on a quality is simply the average of its
questions' grades, and its overall (composite) score is the average across the
five qualities. The whole scoring is then run a second time with the other judge,
so every pipeline ends up with a score from the local \texttt{gemma3:12b} and a
score from the different-family \texttt{gpt-oss-120b}, and the two can be compared.

\subsection{Our results}
\label{sec:res-ours}

Table~\ref{tab:eqg-overall} reports our four pipelines on the five qualities, and
Table~\ref{tab:eqg-bydiff} shows how the overall score changes with difficulty.
\emph{The judging run is still in progress; every cell shown as ``---'' is
pending, and no pending value has been filled in with an estimate.} Once the run
finishes, these tables are completed directly from the measured scores, and the
discussion below is written against the real numbers --- in particular, whether
our pipelines show the same pattern EQGBench found (the first four qualities
saturating while competence guidance separates them), and whether the agentic
pipeline's extra planning and checking earns it any advantage over plain, chain-of-thought
and program-of-thought prompting, especially on harder questions.

\begin{table}[H]
\centering
\caption{Our pipelines on the five EQGBench qualities (each out of $2$; composite
is the average of the five), judged three times per question, reported separately
for the two judges. Pending completion of the judging run; no value is estimated.}
\label{tab:eqg-overall}
\small
\begin{tabular}{|l|l|c|c|c|c|c|c|}
\hline
\textbf{Pipeline} & \textbf{Judge} & \textbf{KP} & \textbf{QT} & \textbf{QQ} &
\textbf{SQ} & \textbf{CG} & \textbf{Composite} \\
\hline
Agentic     & \texttt{gemma3:12b}   & --- & --- & --- & --- & --- & --- \\
            & \texttt{gpt-oss-120b} & --- & --- & --- & --- & --- & --- \\
\hline
Non-agentic & \texttt{gemma3:12b}   & --- & --- & --- & --- & --- & --- \\
            & \texttt{gpt-oss-120b} & --- & --- & --- & --- & --- & --- \\
\hline
CoT         & \texttt{gemma3:12b}   & --- & --- & --- & --- & --- & --- \\
            & \texttt{gpt-oss-120b} & --- & --- & --- & --- & --- & --- \\
\hline
PoT         & \texttt{gemma3:12b}   & --- & --- & --- & --- & --- & --- \\
            & \texttt{gpt-oss-120b} & --- & --- & --- & --- & --- & --- \\
\hline
\end{tabular}
\end{table}

\begin{table}[H]
\centering
\caption{Our overall (composite) score out of $2$ by difficulty, for each judge.
Pending completion of the judging run; no value is estimated.}
\label{tab:eqg-bydiff}
\small
\begin{tabular}{|l|l|c|c|c|}
\hline
\textbf{Pipeline} & \textbf{Judge} & \textbf{Easy} & \textbf{Medium} & \textbf{Hard} \\
\hline
Agentic     & \texttt{gemma3:12b}   & --- & --- & --- \\
            & \texttt{gpt-oss-120b} & --- & --- & --- \\
\hline
Non-agentic & \texttt{gemma3:12b}   & --- & --- & --- \\
            & \texttt{gpt-oss-120b} & --- & --- & --- \\
\hline
CoT         & \texttt{gemma3:12b}   & --- & --- & --- \\
            & \texttt{gpt-oss-120b} & --- & --- & --- \\
\hline
PoT         & \texttt{gemma3:12b}   & --- & --- & --- \\
            & \texttt{gpt-oss-120b} & --- & --- & --- \\
\hline
\end{tabular}
\end{table}

\section{Summary}
\label{sec:res-summary}

This chapter evaluated four question-generation pipelines --- agentic,
non-agentic, chain-of-thought and program-of-thought --- built on the same
pinned local \texttt{gemma3:4b} generator, using the EQGBench protocol: five
dimensions scored $0$--$2$, three-round voting, judged by two independent models
--- the local \texttt{gemma3:12b} and the different-family \texttt{gpt-oss-120b}.
The evaluation deliberately separates generation from measurement, scores an
already-generated corpus of $583$ questions, and reports per-cell sample sizes
because they are not uniform. The measured dimension
scores and the analysis of which dimensions discriminate between pipelines are
completed once the judging run finishes.
